\documentclass[11pt,a4paper]{article}
\usepackage[margin=2.2cm]{geometry}
\usepackage{amsmath,amssymb,mathtools}
\usepackage[T1]{fontenc}
\usepackage[utf8]{inputenc}
\usepackage{lmodern}
\usepackage{hyperref}
\usepackage{listings}
\lstset{basicstyle=\ttfamily\small,breaklines=true,frame=single}

\title{Symbolic Van Vleck Derivation of Quartic Centrifugal Distortion Constants\\
\large A $\tau$-First Formulation in Terms of the $H_{mn}$ Blocks}
\author{Codex + custom SymPy workflow}
\date{\today}

\begin{document}
\maketitle

\section{Aim}
The purpose of this implementation is to derive symbolically the quartic rotational block
of the rovibrational effective Hamiltonian starting from the standard block decomposition
\begin{equation}
H = H_{20} + H_{02} + H_{12} + H_{30} + H_{22} + H_{40},
\end{equation}
where
\begin{align}
H_{20} &= \sum_i \hbar\omega_i\left(a_i^\dagger a_i + \frac12\right),\\
H_{02} &= \text{rigid-rotor block},\\
H_{12} &= \frac12 \sum_{\alpha\beta k} \mu_{\alpha\beta,k} q_k J_\alpha J_\beta,\\
H_{22} &= \frac14 \sum_{\alpha\beta kl} \mu_{\alpha\beta,kl} q_k q_l J_\alpha J_\beta,\\
H_{30} &= \frac16 \sum_{ijk} \Phi_{ijk} q_i q_j q_k,\\
H_{40} &= \frac1{24} \sum_{ijkl} \Phi_{ijkl} q_i q_j q_k q_l.
\end{align}
The normal coordinates are written as
\begin{equation}
q_i = \sqrt{\frac{\hbar}{2\omega_i}}\,(a_i + a_i^\dagger),
\end{equation}
with bosonic commutator
\begin{equation}
[a_i,a_j^\dagger]=\delta_{ij}.
\end{equation}

The key design choice is that the symbolic derivation stops first at the quartic tensor
\(\tau\), rather than directly at Watson constants. Watson A- or S-reduction constants are
then obtained by linear projection of \(\tau\).

\section{Perturbative Organization}
The perturbation expansion is organized according to the standard rovibrational power counting,
\begin{align}
H^{(0)} &= H_{20} + H_{02},\\
H^{(1)} &= H_{12} + H_{30},\\
H^{(2)} &= H_{22} + H_{40}.
\end{align}
With this assignment, the quartic rotational block of the effective Hamiltonian is naturally
resolved into the following contribution classes:
\begin{itemize}
\item \((H_{12},H_{12})\),
\item \((H_{22})\),
\item \((H_{12},H_{30})\),
\item \((H_{30},H_{30})\),
\item \((H_{40})\).
\end{itemize}

This classification is also used as an explicit filtering criterion in the code, so that only
physically relevant terms are retained during the symbolic expansion.

\section{Implementation}
The implementation is contained in
\begin{center}
\texttt{/Users/vincenzobarone/centrifugal/derive\_watson\_quartic\_vanvleck.py}
\end{center}

\subsection{Operator Algebra}
The code represents each operator monomial explicitly as
\begin{itemize}
\item a vibrational word, i.e. a tuple of \texttt{("a",k)} and \texttt{("ad",k)},
\item a non-commuting rotational word in \(J_x,J_y,J_z\),
\item metadata tracking the \(H_{mn}\) origin of the term.
\end{itemize}

The following operations are implemented:
\begin{itemize}
\item multiplication of operator expressions,
\item commutators \([A,B]=AB-BA\),
\item recursive bosonic normal ordering using
\[
a_i a_j^\dagger \rightarrow a_j^\dagger a_i + \delta_{ij}.
\]
\end{itemize}

\subsection{Van Vleck Construction}
The effective Hamiltonian is constructed by a Lie transform,
\begin{equation}
K = e^{S} H e^{-S}
  = H + [S,H] + \frac{1}{2!}[S,[S,H]] + \frac{1}{3!}[S,[S,[S,H]]] + \cdots,
\end{equation}
with
\begin{equation}
S = \sum_{n\ge1} S^{(n)}.
\end{equation}
At each perturbative order \(m\), \(S^{(m)}\) is chosen so as to remove the vibrationally
off-diagonal part of \(K^{(m)}\).

\subsection{Filtering Layer}
To keep the symbolic derivation tractable, each operator monomial is filtered using three
criteria:
\begin{itemize}
\item rotational degree \(\le 4\),
\item even bosonic degree,
\item allowed contribution class in the set
\[
(H_{12},H_{12}),\ (H_{22}),\ (H_{12},H_{30}),\ (H_{30},H_{30}),\ (H_{40}).
\]
\end{itemize}
Only surviving terms are passed to the full normal-ordering and Lie-transform machinery.

\subsection{Physical Output}
After bosonic normal ordering and projection onto the vibrational ground state,
only the quartic rotational block is retained. The primary physical output is the symmetric
quartic tensor
\begin{equation}
(\tau_{xxxx},\tau_{yyyy},\tau_{zzzz},\tau_{xxyy},\tau_{xxzz},\tau_{yyzz}),
\end{equation}
which multiplies
\begin{equation}
J_x^4,\ J_y^4,\ J_z^4,\ J_x^2J_y^2,\ J_x^2J_z^2,\ J_y^2J_z^2.
\end{equation}
Watson constants in either A or S reduction are obtained afterwards by linear projection.

\section{Commands Used}
\begin{lstlisting}
python3 derive_watson_quartic_vanvleck.py --max-order 2
python3 derive_watson_quartic_vanvleck.py --max-order 3
python3 derive_watson_quartic_vanvleck.py --max-order 4
\end{lstlisting}

\section{Notation}
For the reduced symbolic runs discussed below (one mode, diagonal ansatz), the following
publication-style notation is used:
\begin{align}
\mu_{xx}^{(1)} &\equiv \mu_{xx,1}, &
\mu_{yy}^{(1)} &\equiv \mu_{yy,1}, &
\mu_{zz}^{(1)} &\equiv \mu_{zz,1},\\
\mu_{xx}^{(2)} &\equiv \mu_{xx,11}, &
\mu_{yy}^{(2)} &\equiv \mu_{yy,11}, &
\mu_{zz}^{(2)} &\equiv \mu_{zz,11},\\
\Phi^{(3)} &\equiv \Phi_{111}, &
\Phi^{(4)} &\equiv \Phi_{1111}, &
\omega &\equiv \omega_1.
\end{align}

\section{Results}
\subsection{Second Order}
\paragraph{Contribution class \((H_{12},H_{12})\).}
At second order, the quartic tensor receives contributions only from the class
\((H_{12},H_{12})\). In index form, this means that the effective quartic tensor is bilinear
in the first rovibrational coupling tensor \(\mu_{\alpha\beta,k}\mu_{\gamma\delta,k}\), with
the expected harmonic denominators. In the monomial basis used by the code,
\[
J_x^4,\ J_y^4,\ J_z^4,\ J_x^2J_y^2,\ J_x^2J_z^2,\ J_y^2J_z^2,
\]
the exact second-order coefficients are
\begin{align}
\tau_{\alpha\alpha\alpha\alpha}^{(2)} &=
-\sum_k \frac{\mu_{\alpha\alpha,k}^2}{8\omega_k^2},\\
\tau_{\alpha\alpha\beta\beta}^{(2)} &=
-\sum_k \frac{\mu_{\alpha\alpha,k}\mu_{\beta\beta,k}}{4\omega_k^2},
\qquad \alpha \neq \beta.
\end{align}
This is the standard known second-order result in the present normalization, and the symbolic
code is used here as an internal validation benchmark: it must reproduce exactly this structure
before the higher-order channels are analyzed.

In the reduced one-mode diagonal run used for explicit symbolic verification, this becomes
\begin{align}
\tau_{xxxx} &= -\frac{\left(\mu_{xx}^{(1)}\right)^2}{8\omega^2}, &
\tau_{yyyy} &= -\frac{\left(\mu_{yy}^{(1)}\right)^2}{8\omega^2}, &
\tau_{zzzz} &= -\frac{\left(\mu_{zz}^{(1)}\right)^2}{8\omega^2},\\
\tau_{xxyy} &= -\frac{\mu_{xx}^{(1)}\mu_{yy}^{(1)}}{4\omega^2}, &
\tau_{xxzz} &= -\frac{\mu_{xx}^{(1)}\mu_{zz}^{(1)}}{4\omega^2}, &
\tau_{yyzz} &= -\frac{\mu_{yy}^{(1)}\mu_{zz}^{(1)}}{4\omega^2}.
\end{align}

Projection of this tensor onto Watson A reduction yields, for the same reduced run,
\begin{align}
D_J &= \frac{\mu_{xx}^{(1)}\mu_{yy}^{(1)}}{8\omega^2},\\
D_{JK} &= \frac{-2\mu_{xx}^{(1)}\mu_{yy}^{(1)} + \mu_{xx}^{(1)}\mu_{zz}^{(1)} + \mu_{yy}^{(1)}\mu_{zz}^{(1)}}{8\omega^2},\\
D_K &= \frac{\mu_{xx}^{(1)}\mu_{yy}^{(1)} - \mu_{xx}^{(1)}\mu_{zz}^{(1)} - \mu_{yy}^{(1)}\mu_{zz}^{(1)} + \left(\mu_{zz}^{(1)}\right)^2}{8\omega^2},\\
d_1 &= \frac{\mu_{zz}^{(1)}\left(\mu_{xx}^{(1)}-\mu_{yy}^{(1)}\right)}{8\omega^2},\\
d_2 &= \frac{\left(\mu_{xx}^{(1)}\right)^2 - \left(\mu_{yy}^{(1)}\right)^2 - 2\mu_{zz}^{(1)}\left(\mu_{xx}^{(1)}-\mu_{yy}^{(1)}\right)}{16\omega^2}.
\end{align}
Thus, at second order, the quartic Watson constants arise exclusively from the
\((H_{12},H_{12})\) channel. No contribution from \(H_{22}\), \(H_{30}\), or \(H_{40}\)
appears yet.

\subsection{Third Order}
In the current symbolic run, no new quartic ground-state contribution is produced at third order:
\(K^{(3)}\) does not add new \(J^4\) terms. In terms of the \(H_{mn}\) classification, the
candidate odd-order channels do not generate an additional effective contribution to \(\tau\)
in the present setup.

\subsection{Fourth Order}
\paragraph{Contribution classes \((H_{22})\), \((H_{12},H_{30})\), \((H_{30},H_{30})\), \((H_{40})\).}
At fourth order, the code produces nonzero corrections to the quartic tensor. In general one
expects contributions from the classes
\begin{itemize}
\item terms linear in \(\mu_{\alpha\alpha}^{(2)}\), associated with \(H_{22}\),
\item terms linear in \(\Phi^{(4)}\), associated with \(H_{40}\),
\item terms quadratic in \(\Phi^{(3)}\), associated with \((H_{30},H_{30})\),
\item mixed terms \(\mu_{\alpha\alpha}^{(1)}\mu_{\beta\beta}^{(2)}\Phi^{(3)}\), associated with the \((H_{12},H_{30})\) channel.
\end{itemize}

Accordingly, the fourth-order quartic block is expected to have the schematic structure
\begin{align}
\tau^{(4)} &= \tau[H_{22}] + \tau[H_{40}] + \tau[H_{30},H_{30}] + \tau[H_{12},H_{30}],
\end{align}
with schematic dependence
\begin{align}
\tau[H_{22}] &\sim \hbar\,\mu^{(2)},\\
\tau[H_{40}] &\sim \hbar\,\mu^{(1)}\mu^{(1)}\Phi^{(4)},\\
\tau[H_{30},H_{30}] &\sim \hbar\,\mu^{(1)}\mu^{(1)}\left(\Phi^{(3)}\right)^2,\\
\tau[H_{12},H_{30}] &\sim \hbar\,\mu^{(1)}\mu^{(2)}\Phi^{(3)}.
\end{align}

In the present one-mode run, however, the explicit symbolic output separates only into the two
nonzero channels
\begin{align}
\tau^{(4)} = \tau[H_{30},H_{30}] + \tau[H_{12},H_{30}],
\end{align}
while no standalone \(\tau[H_{22}]\) or \(\tau[H_{40}]\) term survives in this restricted setup.

For the diagonal tensor elements,
\begin{align}
\tau_{\alpha\alpha\alpha\alpha}^{(4)} &=
-\frac{\hbar\,\mu_{\alpha\alpha}^{(1)}\mu_{\alpha\alpha}^{(1)}\left(\Phi^{(3)}\right)^2}{1728\,\omega^7}
+\frac{\hbar\,\mu_{\alpha\alpha}^{(1)}\mu_{\alpha\alpha}^{(2)}\Phi^{(3)}}{144\,\omega^5}.
\end{align}
The mixed components \(\tau_{\alpha\alpha\beta\beta}^{(4)}\) have the same two-channel structure.
Explicit one-mode formulas are reported in Appendix~\ref{app:onemode}.

\section{Scope and Generality}
The explicit expressions above were obtained in a reduced symbolic setup (one mode, diagonal
tensor ansatz) to keep the closed-form formulas readable. The implementation itself is more general:
\begin{itemize}
\item rotational operators remain non-commuting throughout the derivation,
\item bosonic normal ordering is exact,
\item diagonal and off-diagonal parts are separated recursively to construct \(S^{(n)}\),
\item the filtering layer keeps the multi-mode derivation tractable.
\end{itemize}

Explicit tensor symmetries such as
\(\Phi_{ijk}=\Phi_{\mathrm{perm}}\) and \(\Phi_{ijkl}=\Phi_{\mathrm{perm}}\)
may be imposed afterwards by symbolic substitutions if a more reduced final form is desired.

\appendix
\section{Explicit One-Mode Fourth-Order \texorpdfstring{$\tau$}{tau} Formulas}
\label{app:onemode}
For completeness, the explicit one-mode fourth-order tensor coefficients obtained in the reduced
symbolic run are
\begin{align}
\tau_{xxxx}^{(4)} &=
-\frac{\hbar\,\left(\mu_{xx}^{(1)}\right)^2\left(\Phi^{(3)}\right)^2}{1728\,\omega^7}
+\frac{\hbar\,\mu_{xx}^{(1)}\mu_{xx}^{(2)}\Phi^{(3)}}{144\,\omega^5},\\
\tau_{yyyy}^{(4)} &=
-\frac{\hbar\,\left(\mu_{yy}^{(1)}\right)^2\left(\Phi^{(3)}\right)^2}{1728\,\omega^7}
+\frac{\hbar\,\mu_{yy}^{(1)}\mu_{yy}^{(2)}\Phi^{(3)}}{144\,\omega^5},\\
\tau_{zzzz}^{(4)} &=
-\frac{\hbar\,\left(\mu_{zz}^{(1)}\right)^2\left(\Phi^{(3)}\right)^2}{1728\,\omega^7}
+\frac{\hbar\,\mu_{zz}^{(1)}\mu_{zz}^{(2)}\Phi^{(3)}}{144\,\omega^5},\\
\tau_{xxyy}^{(4)} &=
-\frac{\hbar\,\left(2\mu_{xx}^{(1)}\mu_{yy}^{(1)}+\left(\mu_{xy}^{(1)}\right)^2+2\mu_{xy}^{(1)}\mu_{yx}^{(1)}+\left(\mu_{yx}^{(1)}\right)^2\right)\left(\Phi^{(3)}\right)^2}{1728\,\omega^7}\nonumber\\
&\quad +\frac{\hbar\,\Phi^{(3)}\left(\mu_{xx}^{(1)}\mu_{yy}^{(2)}+\mu_{xy}^{(1)}\mu_{xy}^{(2)}+\mu_{xy}^{(1)}\mu_{yx}^{(2)}+\mu_{yx}^{(1)}\mu_{xy}^{(2)}+\mu_{yx}^{(1)}\mu_{yx}^{(2)}+\mu_{yy}^{(1)}\mu_{xx}^{(2)}\right)}{144\,\omega^5},\\
\tau_{xxzz}^{(4)} &=
-\frac{\hbar\,\left(2\mu_{xx}^{(1)}\mu_{zz}^{(1)}+\left(\mu_{xz}^{(1)}\right)^2+2\mu_{xz}^{(1)}\mu_{zx}^{(1)}+\left(\mu_{zx}^{(1)}\right)^2\right)\left(\Phi^{(3)}\right)^2}{1728\,\omega^7}\nonumber\\
&\quad +\frac{\hbar\,\Phi^{(3)}\left(\mu_{xx}^{(1)}\mu_{zz}^{(2)}+\mu_{xz}^{(1)}\mu_{xz}^{(2)}+\mu_{xz}^{(1)}\mu_{zx}^{(2)}+\mu_{zx}^{(1)}\mu_{xz}^{(2)}+\mu_{zx}^{(1)}\mu_{zx}^{(2)}+\mu_{zz}^{(1)}\mu_{xx}^{(2)}\right)}{144\,\omega^5},\\
\tau_{yyzz}^{(4)} &=
-\frac{\hbar\,\left(2\mu_{yy}^{(1)}\mu_{zz}^{(1)}+\left(\mu_{yz}^{(1)}\right)^2+2\mu_{yz}^{(1)}\mu_{zy}^{(1)}+\left(\mu_{zy}^{(1)}\right)^2\right)\left(\Phi^{(3)}\right)^2}{1728\,\omega^7}\nonumber\\
&\quad +\frac{\hbar\,\Phi^{(3)}\left(\mu_{yy}^{(1)}\mu_{zz}^{(2)}+\mu_{yz}^{(1)}\mu_{yz}^{(2)}+\mu_{yz}^{(1)}\mu_{zy}^{(2)}+\mu_{zy}^{(1)}\mu_{yz}^{(2)}+\mu_{zy}^{(1)}\mu_{zy}^{(2)}+\mu_{zz}^{(1)}\mu_{yy}^{(2)}\right)}{144\,\omega^5}.
\end{align}

\end{document}
